\documentclass[acmsmall,screen,nonacm,11pt,natbib=false]{acmart}
\geometry{letterpaper,left=1in,right=1in,top=1in,bottom=1in}
\settopmatter{printacmref=false,printccs=false,printfolios=true}
\setcopyright{none}
\renewcommand{\bibliofont}{\normalsize}
\newcommand{\OPT}{\operatorname{OPT}}
\newcommand{\F}{F}
\newcommand{\emptysetv}{\varnothing}
\title{Optimal Deterministic Incremental--Decremental Submodular Maximization}
\author{Zijian Zeng}
\affiliation{\institution{UCSI University}\country{Malaysia}}
\email{1002266693@ucsiuniversity.edu.my}
\renewcommand{\shortauthors}{Zijian Zeng}
\date{}
\begin{document}
\begin{abstract}
\normalsize
Incremental--decremental maximization asks for one ordering of a finite ground
set such that every prefix has high utility. For two nonnegative monotone
submodular functions $h$ and $g$, the utility of a prefix $S$ is
$h(S)+g(E\setminus S)$, and its benchmark is the best set of the same size.
Disser, Klimm, Lutz and Strubberg established a deterministic competitive
upper bound of $1+e/(e-1)$ and a structural lower bound of $2$.

We give a deterministic $2$-competitive algorithm for the full class,
including arbitrary unequal endpoint totals. It uses $O(|E|^2)$ exact
value-oracle queries and polynomial additional computation. If $h(E)\ge g(E)$,
the algorithm greedily maximizes the actual objective marginal at each step,
continuing through negative gains. Otherwise it swaps the functions and
reverses the resulting order. The proof retains the initial objective as a
budget in a submodular union comparison; endpoint orientation controls the
negative-gain suffix. Combined with the existing lower bound, which we
reproduce with attribution, this determines the optimal uniform deterministic
constant as $2$, both with unrestricted computation and in the polynomial
exact-value-oracle model.
\end{abstract}
\maketitle
\clearpage

\section{The problem and the sharp guarantee}

An infrastructure renewal order must maintain useful service throughout a
transition: some elements already operate under a new system, while the
others remain in the old system. The incremental--decremental model of
Disser, Klimm, Lutz and Strubberg~\cite{DKLS26} captures this requirement
by evaluating both parts of a partition. The mathematical difficulty is to
choose a \emph{single} order that performs well at every cardinality,
rather than an unrelated approximate solution for each cardinality.

Let $E$ be a finite ground set of size $n$. A set function $q:2^E\to
\mathbb R_{\ge0}$ is normalized if $q(\varnothing)=0$, monotone if
$q(A)\le q(B)$ whenever $A\subseteq B$, and submodular if
\[
 q(A)+q(B)\ge q(A\cup B)+q(A\cap B)\qquad(A,B\subseteq E).
\]
Initially assume that $h$ and $g$ have all three properties. Define
\begin{equation}\label{eq:objective}
 F(S)=h(S)+g(E\setminus S),\qquad
 \OPT_k=\max_{\substack{O\subseteq E\\|O|=k}}F(O).
\end{equation}
For a permutation $\pi=(e_1,\ldots,e_n)$, write
$S_k=\{e_1,\ldots,e_k\}$ and $S_0=\varnothing$.
The permutation is $r$-competitive if
\begin{equation}\label{eq:competitive}
 rF(S_k)\ge\OPT_k\qquad(0\le k\le n).
\end{equation}
This inequality convention includes zero objectives without dividing by zero.
The endpoint cardinalities each have a unique feasible set.

The previous full-class bounds are $2$ and $1+e/(e-1)\approx2.582$:
Corollary~2 and Proposition~11 of~\cite{DKLS26} give the upper and lower
bounds, respectively, and its Section~5 asks to close the gap. Its
Proposition~17 already gives a factor of $2$ when $h(E)=g(E)$.
We address the entire pair class, with no equality requirement on these totals.

\begin{theorem}\label{thm:main}
For every finite $E$ and normalized nonnegative monotone submodular
$h,g:2^E\to\mathbb R_{\ge0}$, a deterministic algorithm returns one permutation
satisfying
\[
 2F(S_k)\ge\OPT_k\qquad\text{simultaneously for every }k=0,\ldots,n.
\]
It uses $O(n^2)$ exact value-oracle queries and polynomial additional
computation in a unit-cost exact-arithmetic and comparison model.
The optimal uniform deterministic competitive constant is $2$ for both
unrestricted algorithms and polynomial-oracle algorithms.
The same constants hold for nonnegative monotone submodular inputs without
normalization.
\end{theorem}

The algorithmic upper bound is the contribution here. The impossibility of
any uniform constant below $2$ is prior work~\cite[Proposition~11]{DKLS26};
Section~\ref{sec:lower} reproduces a coverage-family proof for completeness.
Thus the result determines the exact worst-case cost of maintaining a single
nested sequence for this model, with an efficient construction.

The key distinction from ordinary monotone submodular greedy analysis is
that $F$ may decrease. We preserve its positive initial value in the proof
rather than subtracting it. Before greedy gains become negative, that initial
value pays for the loss in comparing an optimal set with its union with a
prefix. Afterward, orienting the order toward the larger endpoint supplies
the guarantee directly. Both parts use elementary diminishing returns; no
claim of a more general nonmonotone greedy theorem is needed.

\paragraph{Relation to nearby results.}
The double-greedy algorithm in~\cite{DKLS26} compares separate $h$ and $g$
marginals when extending opposite ends of an order. Its exact worst-case
ratio is $1+e/(e-1)$, even with modular $g$
\cite[Proposition~19 and Corollary~20]{DKLS26}. Our algorithm instead commits
to one direction and maximizes the marginal of the \emph{combined} objective.
Consequently this is a different algorithm, rather than a sharper analysis
of that double-greedy rule.

For monotone incremental maximization, the subcase $g=0$,
\cite{BBBCDL26} reports an upper bound of approximately $1.373$ and a lower bound of $1.25$.
Its improved upper bound sacrifices computational efficiency and does not
supply an upper bound for arbitrary pairs $h,g$. The factor $2$ proved here
is sharp over the pair class; it is not asserted to be sharp for that subcase.
For another nearby comparison, Chen and Kuhnle~\cite{CK26} analyze
curvature-dependent greedy with pruning under a fixed at-most-$k$ constraint.
Pruning can remove selected elements, so that result does not itself provide
one permutation at every exact cardinality. These comparisons concern the
stated models and theorems, rather than all possible consequences of those
papers. A bounded primary-literature check found no equivalent full-class
guarantee as of November~1, 2026; this does not establish universal priority.

\section{The algorithm}\label{sec:algorithm}

Fix a total order on $E$ for deterministic tie breaking. For a function $W$,
write $W(e\mid P)=W(P\cup\{e\})-W(P)$ when $e\notin P$.
Let $H=h(E)$ and $G=g(E)$.

\begin{enumerate}
\item If $H\ge G$, start with $P_0=\varnothing$. For $i=1,\ldots,n$ choose
\[
 e_i\in\arg\max_{e\in E\setminus P_{i-1}}F(e\mid P_{i-1}),
 \qquad P_i=P_{i-1}\cup\{e_i\},
\]
breaking ties by the fixed order. Return $(e_1,\ldots,e_n)$.
\item If $H<G$, apply the preceding greedy construction to the swapped
objective $F'(A)=g(A)+h(E\setminus A)$, and return its order in reverse.
\end{enumerate}

The construction always continues until every element has been selected,
including when the maximum available marginal is negative. It does not
query $\OPT_k$. For $E=\varnothing$ it returns the empty permutation.
The endpoint comparison is part of the algorithm, and will be used in the
negative-gain case of the proof.

There are $n(n+1)/2$ candidate evaluations. Each combined objective value
uses two summand-oracle queries, and caching the current objective permits
each marginal comparison with constant additional scalar arithmetic.
Together with the endpoint queries this is at most $n(n+1)+2$ summand
queries for normalized inputs. Set construction and bookkeeping require
only polynomial additional computation. The model permits exact arithmetic
and comparisons on oracle answers at unit cost; no bit-complexity guarantee
for arbitrarily encoded real answers is asserted.

\section{Proof of the simultaneous upper bound}\label{sec:proof}

We first consider the orientation $H\ge G$. Complementation preserves
submodularity: apply the defining inequality for $g$ to $E\setminus A$
and $E\setminus B$, and interchange union and intersection. Thus
$S\mapsto g(E\setminus S)$, and hence $F$, is submodular. Although $F$
need not be monotone, its decomposition yields two useful bounds.

\begin{lemma}\label{lem:budget}
For every $O,P\subseteq E$,
\begin{align}
 F(O)&\le H+G,\label{eq:global}\\
 F(O)&\le F(O\cup P)+G.\label{eq:union}
\end{align}
\end{lemma}
\begin{proof}
Monotonicity bounds $h(O)$ by $H$ and $g(E\setminus O)$ by $G$,
proving~\eqref{eq:global}. Moreover,
\begin{align*}
 F(O)-F(O\cup P)
 &=h(O)-h(O\cup P)\\
 &\quad+g(E\setminus O)-g(E\setminus(O\cup P))\le G.
\end{align*}
The first difference is nonpositive by monotonicity of $h$; the second
is at most $G$ by monotonicity and nonnegativity of $g$.
\end{proof}

Let $a_i=F(P_i)-F(P_{i-1})$ be the gains of the greedy order. These gains
are nonincreasing, including after they become negative. Indeed, for
$e\notin P_i$, diminishing returns and the preceding maximal choice give
\[
 F(e\mid P_i)\le F(e\mid P_{i-1})\le a_i.
\]
Maximizing the left side over the remaining elements proves
$a_{i+1}\le a_i$ for $i<n$. Also $F(\varnothing)=G$ and $F(E)=H$.

Fix any internal cardinality $1\le k<n$, put $P=P_k$ and $d=a_{k+1}$,
and let $O$ be any set of size $k$.

\paragraph{Case 1: $d\ge0$.}
Adding the elements of $O\setminus P$ one by one, diminishing returns gives
\begin{equation}\label{eq:extend}
 F(O\cup P)-F(P)
 \le\sum_{e\in O\setminus P}F(e\mid P)
 \le |O\setminus P|d\le kd.
\end{equation}
The middle inequality follows because $d$ is the maximum marginal at $P$.
The last inequality uses $d\ge0$ and $|O\setminus P|\le k$.
Since the gains are nonincreasing,
\begin{equation}\label{eq:charge}
 kd\le\sum_{i=1}^k a_i=F(P)-G.
\end{equation}
Combining~\eqref{eq:union}, \eqref{eq:extend} and~\eqref{eq:charge} yields
\[
 F(O)\le F(P)+G+kd\le 2F(P).
\]
The initial value $G$ in~\eqref{eq:charge} cancels exactly the union-loss
budget in~\eqref{eq:union}. The argument allows $d=0$ and arbitrary
intersection between $O$ and $P$.

\paragraph{Case 2: $d<0$.}
Every remaining gain is at most $d<0$, so
\[
 H=F(E)=F(P)+\sum_{i=k+1}^{n}a_i\le F(P).
\]
Using~\eqref{eq:global} and the chosen orientation $H\ge G$ gives
\[
 F(O)\le H+G\le2H\le2F(P).
\]

The set $O$ was arbitrary. Thus this same order is $2$-competitive at every
internal cardinality. At $k=0,n$ the unique feasible set is itself optimal,
so the guarantee also holds there. Nonnegativity handles even zero endpoint
values, and the proof never divides by a function value. There is no internal
cardinality to consider for $n\le1$.

\paragraph{Reversing the orientation.}
If $H<G$, swap $h$ and $g$ and let $(u_1,\ldots,u_n)$ be the greedy order
for $F'$. Its totals satisfy $H'=G\ge G'=H$, so the argument just proved
applies. Write $Q_j=\{u_1,\ldots,u_j\}$. A size-$k$ prefix $S_k$ of
the returned reverse order is $E\setminus Q_{n-k}$. Therefore
\[
 F(S_k)=F'(Q_{n-k}),\qquad
 \OPT_k=\max_{|A|=n-k}F'(A).
\]
The second equality uses the bijection $A=E\setminus O$ between all
size-$k$ comparators and size-$(n-k)$ comparators. The swapped guarantee
now gives $2F(S_k)\ge\OPT_k$ simultaneously for every $k$.
This completes the normalized upper-bound proof.

\paragraph{Nonnegative baselines.}
For unnormalized nonnegative monotone inputs set
$\bar h(S)=h(S)-h(\varnothing)$,
$\bar g(S)=g(S)-g(\varnothing)$, and
$B=h(\varnothing)+g(\varnothing)\ge0$.
The barred functions satisfy the normalized hypotheses. Apply the algorithm,
including its endpoint comparison, to these functions. With
$\bar F=\bar h+\bar g\circ(E\setminus\cdot)$, we have
$F=\bar F+B$ and $\OPT_k=\overline{\OPT}_k+B$. Hence
\[
 2F(S_k)=2\bar F(S_k)+2B
 \ge\overline{\OPT}_k+2B\ge\OPT_k.
\]
This costs two extra baseline queries. It is the usual normalization transfer,
also used in~\cite[Eq.~(5)]{DKLS26}; negative baselines are outside the model.

\section{The matching lower bound from prior work}\label{sec:lower}

For completeness we reproduce the curvature-one coverage construction of
Disser, Klimm, Lutz and Strubberg~\cite[Proposition~11]{DKLS26}.
For $n\ge2$ let $E=\{a,b_1,\ldots,b_{n-1}\}$ and define
\[
 q(S)=\begin{cases}
 n-1,&a\in S,\\
 |S|,&a\notin S.
 \end{cases}
\]
This is a normalized nonnegative monotone coverage function: there are $n-1$
unit-weight features, $a$ covers all of them, and $b_i$ covers only feature
$i$. Set $h=g=q$. Each summand is submodular.

At cardinalities $1$ and $n-1$ the optimum is $2(n-1)$, attained by
$\{a\}$ and $E\setminus\{a\}$, respectively. This value is also an upper
bound because each summand is at most $n-1$. In any permutation, if $a$
is not first, its first prefix is a singleton $b_i$ and has objective value
$1+(n-1)=n$. If $a$ is first, its prefix of size $n-1$ contains $a$, and
the complement is one $b_i$; again its objective value is $n$.
Every permutation therefore has competitive ratio at least
\[
 \frac{2(n-1)}{n}=2-\frac2n.
\]
For every $r<2$, choose an integer $n>2/(2-r)$. This instance has no
$r$-competitive permutation, even with unlimited computation and full
knowledge of both functions.

Formally, let $C_{\mathrm{struct}}$ be the infimum of uniform competitive
constants of deterministic algorithms without computational restrictions,
and let $C_{\mathrm{poly}}$ be the same infimum restricted to the exact
polynomial-oracle model. The lower family and the upper algorithm give
\[
 2\le C_{\mathrm{struct}}\le C_{\mathrm{poly}}\le2.
\]
Thus both equal $2$, completing Theorem~\ref{thm:main}. No assumption that
a finite lower-bound instance attains the limiting value is needed.
Normalized instances already force this lower bound, so allowing nonnegative
unnormalized baselines does not alter either constant.

\section{Scope and research methodology}\label{sec:limits}

The guarantee is deterministic and applies to one permutation for every
cardinality, arbitrary unequal totals, and arbitrary exact ties resolved by
a fixed rule. Its proof uses the pair decomposition in
Lemma~\ref{lem:budget}; it does not give a uniform guarantee for all
nonmonotone submodular functions. Indeed,~\cite[Proposition~18]{DKLS26}
shows an unbounded deterministic incremental ratio in that larger class.
The result also makes no optimality claim for narrower function classes,
randomized algorithms, or additional feasibility constraints.

Substantial AI assistance was used in developing the algorithm and proof,
searching for counterexamples, implementing finite exact-arithmetic and
numerical diagnostic checks, comparing primary literature, and drafting
this text. Counterexamples to intermediate proof proposals were used to
redirect the analysis; they are not disproofs of the original ordering
problem. The universal result rests on the inequalities in
Section~\ref{sec:proof}, not on finite experiments or model-generated review.
The proof is elementary, and no broader consequences beyond the stated
incremental--decremental model are claimed.

\bibliographystyle{alphaurl}
\bibliography{references}

\section*{AI methodology statement}
OpenAI Codex was used substantively for algorithm and proof discovery,
counterexample search, computational implementation and checks,
primary-literature comparison, an additional model-based examination of the
argument, and preparation of the entire manuscript. These uses materially
affected all sections. The model-based examination is not external peer
review, and no independent human verification is claimed in this document.
The mathematical claims and references require the sole human author's
responsibility before any submission. Finite checks and AI-generated
assessments are not mathematical certificates.
\end{document}
